\scanpage{102}

\uline{Car. \unsure{quelc.}}

\[
\deleted{q([x,y])=\varphi(x,y)^2}
\]

\deleted{Donc si $x+\underline i y\in\Gamma\tilde U$, donc $q(x)=q(y)\inserted{\in\Gamma G_m}$, $\varphi(x,y)=0$, alors $q([x,y])=0$. Si on remplace $x+\underline i y$ par $(a+\underline i b)(x+\underline i y)=\underbrace{ax-by}_{x'}+\underline i(\underbrace{bx+ay}_{y'})$, on trouve $x'\wedge y'=(a^2+b^2)x\wedge y$, $q(x')=q(y')=(a^2+b^2)q(x)$, donc $x'\wedge y'/q(x')=x\wedge y/q(x)$. On trouve}

\[
\deleted{\unsure{x+iy}}\quad x+\underline i y\in\Gamma\tilde U\longrightarrow[x,y]/q(x)\in\Sigma_{-4}
\]

donc

\[
\tilde g:\tilde U\longrightarrow\Sigma_{-4}
\]

passe au quotient en

\[
g:U^{\deleted{\unsure{S[i]}}}\longrightarrow\Sigma_{-4}
\]

(\unsure{indépendam$^{\mathrm t}$|indépendamment} de la car.). En car. 2, on a

\deleted{Supposons $[x,y]/q(x)=[x',y']/q^{2}(x)$ ; donc (quitte à se localiser et à remplacer $x'$ par $\lambda x'$). \unsure{supp.} $q(x)=q(x')$, $[x,y]=[x'^{\deleted{\unsure{+}}},y']$. Supposons qu'on \deleted{a) Supposons $y$} \unsure{soit en car.} \unsure{$\ne2$}}

\[
[x,y]=\varphi(x,y)H\quad\text{avec }\left\{\begin{aligned}&H\text{ base du centre de }\unsure{\tilde{\mathfrak g}|\tilde g}\\ &q(H)=1\end{aligned}\right.
\]

donc $q([x,y])=\varphi(x,y)^2$

Si $\varphi(x,y)=0$ (p.ex. $x+\underline i y\in\Gamma\tilde X^{S[i]}\supset\Gamma(\tilde U)$\unsure{\ldots}) on a $[x,y]=0$. Donc en car. 2 l'application $g$ se spécialise en l'application constante de valeur 0 : rien à faire\ldots !

% archive sheet number: 101
